\section{Introduction}
\label{sec:introduction}

Information is a basic organizing principle of artificial intelligence.
Learning systems compress observations into representations, decision systems
reduce uncertainty about actions, and graph-based systems try to expose
structure that is not immediately visible from local relations alone.
Much of classical information theory, however, is local in flavour: it
quantifies uncertainty in symbols, variables, or pairwise dependencies.
\textbf{Structural Entropy (SE)}~\cite{Li2016} asks a different question.
Instead of only measuring uncertainty at the level of isolated variables, it
studies the uncertainty embedded in the organization of a whole system.

This perspective is especially natural for graphs.
Real graph data are rarely just collections of edges.
Social networks contain groups and subgroups, biological systems organize into
functional modules, and modern AI pipelines increasingly generate interaction
graphs whose usefulness depends on whether they admit meaningful hierarchies.
In such settings, one often wants more than a flat partition or a predictive
score.
One wants a structural description of the system: which components belong
together, at what scale, and according to what coding principle.

SE provides one answer to this problem.
In its original form, it connects graph organization to coding trees and
interprets good hierarchical partitions as efficient information-theoretic
descriptions of the graph~\cite{Li2016}.
This gives SE a dual role.
It is a measure, because it quantifies how much uncertainty remains under a
proposed structural decomposition.
It is also an optimization principle, because minimizing it searches for the
decomposition that best exposes the graph's organization.
That combination has made SE appealing not only for community detection, where
it first gained attention, but also for graph pooling, augmentation,
clustering, control, and domain-specific discovery tasks~\cite{wu2022pooling,wu2023sega,zeng2025sidmJMLR,zhang2021supertad,peng2024zjyun,zhu2023HiTIN}.

The field has grown quickly since the foundational work of Li and Pan in 2016.
What began as a graph-theoretic and coding-theoretic framework for uncovering
community structure now includes variants for directed graphs, dynamic
settings, differentiable deep models, and decision-making
systems~\cite{Li2016,pan2021information,sun2024lsenet,zeng2025sidmJMLR}.
SE-based ideas have appeared in graph neural networks, reinforcement learning,
social-bot analysis, genomic data analysis, traffic modelling, and speech or
vision applications~\cite{yang2024sebot,chen2022pyseat,zou2024transformer,wang2024secodecSpeech}.
At the same time, the literature has become fragmented.
Different papers use SE as a discrete objective, a regularizer, a hierarchy
prior, a robustness signal, or a task-specific structural score.
As a result, readers can now find many applications of the idea, but it is
harder to see the shape of the field as a whole.

This survey reviews the first decade of SE research and organizes it around a
small number of recurring questions.
What exactly does SE measure, and how does it differ from related graph
objectives such as modularity, the map equation, spectral criteria, or other
entropy-based quantities?
What algorithmic families have been developed to minimize or approximate it?
How has the idea been adapted to differentiable graph learning and
reinforcement learning?
Which application areas have produced the most convincing evidence of value,
and how should those claims be evaluated?
By bringing these threads together, we aim to provide both a map of the
literature and a clearer understanding of what SE contributes as an
information-theoretic view of graph intelligence.

Our perspective in this journal version is shaped by two developments.
First, SE research is now large enough to support a genuine survey rather than
a short position piece.
Second, graph intelligence has changed substantially in the same period.
Contemporary systems increasingly combine graph structure with foundation
models, retrieval pipelines, agent memories, and multimodal interaction logs.
These systems create new opportunities for structural abstraction, but they
also raise harder questions about interpretability, stability, and
evaluation.
SE is timely in this setting because it offers a language for reasoning about
organization itself, not only prediction.

The present manuscript substantially extends our earlier conference
survey~\cite{su2025ijcai}.
Relative to that version, the journal manuscript expands the theoretical
background, broadens the algorithm catalogue, updates the application review
through 2026, adds a reproducibility benchmark with the \texttt{glass-jax}
companion library, and develops a fuller discussion of open problems.
The aim is not only to collect papers, but to organize a field that now spans
classical graph theory, modern graph learning, and emerging graph-enhanced AI
systems.

The remainder of the survey follows that logic.
Section~\ref{sec:related} positions SE relative to neighboring graph
objectives and learning paradigms.
Section~\ref{sec:taxonomy} introduces the survey taxonomy and formal framing.
Section~\ref{sec:algorithms} reviews major optimization approaches for SE
minimization.
Section~\ref{sec:learning} discusses learning-based formulations, including
graph-learning and reinforcement-learning variants.
Section~\ref{sec:benchmark} examines evaluation practice through our benchmark
and reproducibility lens.
Section~\ref{sec:discussion} and Section~\ref{sec:open} widen the view to
implications, limitations, and future directions, and
Section~\ref{sec:conclusion} concludes.
